\documentclass[a4paper,12pt]{article}

% ------------------------------------------------------------
% Кодировка и русский язык
% ------------------------------------------------------------
\usepackage[T2A]{fontenc}
\usepackage[utf8]{inputenc}
\usepackage[russian]{babel}

% ------------------------------------------------------------
% Математика и типографика
% ------------------------------------------------------------
\usepackage{amsmath,amssymb,amsthm,mathtools}
\usepackage{helvet}
\renewcommand{\familydefault}{\sfdefault}
\usepackage{icomma}
\usepackage{microtype}
\usepackage{setspace}
\usepackage{parskip}

% ------------------------------------------------------------
% Геометрия страницы
% ------------------------------------------------------------
\usepackage[
  left=22mm,
  right=22mm,
  top=22mm,
  bottom=24mm,
  headheight=15pt
]{geometry}

% ------------------------------------------------------------
% Графика
% ------------------------------------------------------------
\usepackage{tikz}
\usepackage{tkz-euclide}
\usepackage{pgfplots}
\pgfplotsset{compat=1.18}

% ------------------------------------------------------------
% Таблицы и списки
% ------------------------------------------------------------
\usepackage{tabularx,booktabs,longtable}
\usepackage{enumitem}

% ------------------------------------------------------------
% Оформление
% ------------------------------------------------------------
\usepackage{xcolor}
\usepackage{titlesec}
\usepackage{fancyhdr}
\usepackage{hyperref}

\definecolor{accent}{HTML}{008080}
\definecolor{darkaccent}{HTML}{004D4D}
\definecolor{textgray}{HTML}{333333}
\definecolor{lightgray}{HTML}{E8EEEE}

\hypersetup{
  unicode=true,
  colorlinks=true,
  linkcolor=darkaccent,
  urlcolor=accent,
  pdftitle={Чек-лист по аннуитетной схеме кредитования},
  pdfauthor={Лёвин Артём Александрович}
}

\titleformat{\section}
  {\Large\bfseries\color{darkaccent}}
  {\thesection.}{0.65em}{}

\titleformat{\subsection}
  {\large\bfseries\color{accent}}
  {\thesubsection.}{0.6em}{}

\titlespacing*{\section}{0pt}{1.5em}{0.7em}
\titlespacing*{\subsection}{0pt}{1.2em}{0.45em}

\setstretch{1.35}
\setlength{\parindent}{0pt}
\setlength{\emergencystretch}{2em}

\pagestyle{fancy}
\fancyhf{}
\fancyhead[L]{\small\color{textgray}Аннуитетная схема кредитования}
\fancyhead[R]{\small\color{textgray}Николь Саркисьянц}
\fancyfoot[C]{\thepage}
\renewcommand{\headrulewidth}{0.4pt}
\renewcommand{\headrule}{%
  \hbox to\headwidth{\color{lightgray}\leaders\hrule height \headrulewidth\hfill}
}

\newlist{checklist}{itemize}{1}
\setlist[checklist]{
  label=\(\square\),
  leftmargin=*,
  itemsep=0.55em,
  topsep=0.35em
}

\newlist{training}{enumerate}{1}
\setlist[training]{
  label=\textbf{\arabic*.},
  leftmargin=*,
  itemsep=1.05em,
  topsep=0.5em
}

\newcommand{\term}[1]{\textbf{\color{darkaccent}#1}}
\newcommand{\important}[1]{\textbf{#1}}
\newcommand{\rub}{\,\text{руб.}}

\begin{document}

% ============================================================
% Титульный лист
% ============================================================

\begin{titlepage}
  \thispagestyle{empty}
  \centering

  \vspace*{2.4cm}

  {\Large\color{accent}\bfseries ЧЕК-ЛИСТ\par}

  \vspace{1.1cm}

  {\Huge\color{darkaccent}\bfseries
  Аннуитетная схема\\[0.25em]
  кредитования\par}

  \vspace{0.8cm}

  {\Large
  Таблица задолженности, платёж,\\
  переплата и процент переплаты\par}

  \vfill

  {\large
  Лёвин Артём Александрович\\[0.35em]
  эксклюзивно для Николь Саркисьянц\par}

  \vspace{1.2cm}

  {\large\today\par}

  \vspace{2.5cm}

  \begin{tikzpicture}[remember picture,overlay]
    \draw[
      accent,
      line width=1.15pt
    ]
    ([xshift=18mm,yshift=20mm]current page.south west)
      .. controls
      ([xshift=55mm,yshift=34mm]current page.south west)
      and
      ([xshift=-68mm,yshift=7mm]current page.south east)
      ..
    ([xshift=-18mm,yshift=20mm]current page.south east);

    \draw[
      darkaccent,
      line width=0.55pt
    ]
    ([xshift=18mm,yshift=15mm]current page.south west)
      .. controls
      ([xshift=75mm,yshift=3mm]current page.south west)
      and
      ([xshift=-52mm,yshift=29mm]current page.south east)
      ..
    ([xshift=-18mm,yshift=15mm]current page.south east);
  \end{tikzpicture}
\end{titlepage}

% ============================================================
% Основной чек-лист
% ============================================================

\section{Смысл аннуитетной схемы}

\begin{checklist}
  \item \term{Аннуитетная схема} означает, что заёмщик в конце каждого
  расчётного периода вносит \important{одинаковый платёж}.

  \item В начале очередного периода имеется некоторая задолженность перед
  банком.

  \item Сначала на текущую задолженность начисляются проценты.

  \item После начисления процентов формируется общая задолженность перед
  очередным платежом.

  \item Затем из общей задолженности вычитается фиксированный платёж.

  \item Остаток после платежа становится начальной задолженностью следующего
  периода.

  \item После последнего платежа задолженность должна стать равной нулю.
\end{checklist}

\[
\boxed{
\text{начальная задолженность}
\longrightarrow
\text{начисление процентов}
\longrightarrow
\text{платёж}
\longrightarrow
\text{остаток долга}
}
\]

\section{Основные обозначения}

\begin{table}[ht]
  \centering
  \caption{Обозначения, используемые в аннуитетной схеме}
  \renewcommand{\arraystretch}{1.4}
  \begin{tabularx}{\linewidth}{@{}cX@{}}
    \toprule
    \textbf{Обозначение} & \textbf{Смысл} \\
    \midrule
    \(S\) &
    первоначальная сумма кредита; \\

    \(P\%\) &
    годовая процентная ставка, указанная в условии; \\

    \(p\) &
    ставка в долях единицы:
    \(\displaystyle p=\frac{P}{100}\); \\

    \(r\) &
    коэффициент увеличения задолженности:
    \(\displaystyle r=1+p\); \\

    \(x\) &
    одинаковый платёж в конце каждого периода; \\

    \(n\) &
    количество расчётных периодов; \\

    \(D_k\) &
    остаток задолженности после \(k\)-го платежа; \\

    \(\text{ОСВ}\) &
    общая сумма выплат банку; \\

    \(\Pi\) &
    переплата по кредиту; \\

    \(\Pi_{\%}\) &
    процент переплаты. \\
    \bottomrule
  \end{tabularx}
\end{table}

\begin{checklist}
  \item До составления таблицы обязательно переведите ставку из процентов
  в доли единицы:
  \[
  p=\frac{P}{100}.
  \]

  \item После этого введите коэффициент
  \[
  r=1+p.
  \]

  \item Не смешивайте \(P\%\) и \(p\): например, при ставке \(12\%\)
  имеем \(p=0{,}12\), а не \(p=12\).
\end{checklist}

\section{Как формируется одна строка таблицы}

Пусть перед очередным начислением процентов задолженность равна \(A\).

\begin{checklist}
  \item Начисленные проценты:
  \[
  Ap.
  \]

  \item Общая задолженность перед платежом:
  \[
  A+Ap=A(1+p)=Ar.
  \]

  \item Остаток после внесения платежа \(x\):
  \[
  Ar-x.
  \]
\end{checklist}

Следовательно, переход от одного периода к следующему описывается правилом
\[
\boxed{A\longmapsto Ar-x}.
\]

\section{Таблица для кредита на три периода}

\begin{table}[ht]
  \centering
  \caption{Изменение задолженности при трёх равных платежах}
  \renewcommand{\arraystretch}{1.55}
  \small
  \begin{tabularx}{\linewidth}{
    @{}
    c
    >{\raggedright\arraybackslash}X
    >{\raggedright\arraybackslash}X
    >{\raggedright\arraybackslash}X
    @{}
  }
    \toprule
    \textbf{Период} &
    \textbf{Долг в начале} &
    \textbf{Долг после начисления процентов} &
    \textbf{Долг после платежа} \\
    \midrule

    1 &
    \(S\) &
    \(Sr\) &
    \(Sr-x\) \\

    2 &
    \(Sr-x\) &
    \((Sr-x)r\) &
    \((Sr-x)r-x\) \\

    3 &
    \((Sr-x)r-x\) &
    \(\bigl((Sr-x)r-x\bigr)r\) &
    \(\bigl((Sr-x)r-x\bigr)r-x\) \\
    \bottomrule
  \end{tabularx}
\end{table}

После последнего платежа кредит должен быть погашен полностью, поэтому
\[
\boxed{\bigl((Sr-x)r-x\bigr)r-x=0}.
\]

\section{Вывод платежа для трёх периодов}

Последовательно раскроем скобки:
\begin{align*}
\bigl((Sr-x)r-x\bigr)r-x
  &=\bigl(Sr^2-xr-x\bigr)r-x\\
  &=Sr^3-xr^2-xr-x.
\end{align*}

Условие полного погашения принимает вид
\[
Sr^3-xr^2-xr-x=0.
\]

Вынесем \(-x\):
\[
Sr^3-x(r^2+r+1)=0.
\]

Перенесём слагаемое с \(x\):
\[
x(r^2+r+1)=Sr^3.
\]

Следовательно,
\[
\boxed{
x=\frac{Sr^3}{r^2+r+1}
}.
\]

\begin{checklist}
  \item Формулу не следует использовать без обоснования.

  \item В полном решении сначала составьте таблицу или рекуррентную модель,
  затем запишите условие полного погашения.

  \item Обратите внимание на закономерность степеней:
  \[
  Sr^3-xr^2-xr-x.
  \]

  \item Степень \(r\) при начальной сумме равна числу периодов, а при
  платежах последовательно уменьшается от \(2\) до \(0\).
\end{checklist}

\section{Общая формула для \(n\) периодов}

Пусть
\[
D_0=S,
\]
а \(D_k\) --- задолженность после \(k\)-го платежа.

Тогда
\[
\boxed{D_k=rD_{k-1}-x}.
\]

Первые значения:
\begin{align*}
D_1&=Sr-x,\\
D_2&=Sr^2-x(r+1),\\
D_3&=Sr^3-x(r^2+r+1).
\end{align*}

В общем случае
\[
\boxed{
D_n=Sr^n-x\left(r^{n-1}+r^{n-2}+\dots+r+1\right)
}.
\]

Так как после последнего платежа \(D_n=0\), получаем
\[
x\left(r^{n-1}+r^{n-2}+\dots+r+1\right)=Sr^n.
\]

Следовательно,
\[
\boxed{
x=
\frac{Sr^n}
{r^{n-1}+r^{n-2}+\dots+r+1}
}.
\]

При \(p>0\), то есть при \(r\ne 1\), сумму геометрической прогрессии можно
записать так:
\[
1+r+r^2+\dots+r^{n-1}
=
\frac{r^n-1}{r-1}
=
\frac{r^n-1}{p}.
\]

Поэтому
\[
\boxed{
x=\frac{Spr^n}{r^n-1}
}.
\]

При нулевой процентной ставке \(p=0\) проценты не начисляются, поэтому
\[
\boxed{x=\frac{S}{n}}.
\]

\section{Четыре периода и факторизация знаменателя}

Для четырёх равных платежей:
\[
\boxed{
x=\frac{Sr^4}{r^3+r^2+r+1}
}.
\]

Знаменатель можно разложить на множители группировкой:
\begin{align*}
r^3+r^2+r+1
  &=r^2(r+1)+(r+1)\\
  &=(r+1)(r^2+1).
\end{align*}

Следовательно,
\[
\boxed{
x=\frac{Sr^4}{(r+1)(r^2+1)}
}.
\]

\begin{checklist}
  \item Если выражение содержит четыре слагаемых, проверьте возможность
  группировки.

  \item Факторизация может существенно упростить дальнейшие вычисления.

  \item После разложения обязательно проверьте результат обратным раскрытием
  скобок.
\end{checklist}

\section{Общая сумма выплат}

Если кредит погашается за \(n\) периодов и каждый платёж равен \(x\), то
\[
\boxed{\text{ОСВ}=nx}.
\]

Для трёх периодов:
\[
\boxed{\text{ОСВ}=3x}.
\]

Для четырёх периодов:
\[
\boxed{\text{ОСВ}=4x}.
\]

\section{Переплата}

\term{Переплата} --- это сумма, выплаченная банку сверх первоначально
полученной суммы кредита.

Наиболее удобная формула:
\[
\boxed{\Pi=\text{ОСВ}-S}.
\]

Так как \(\text{ОСВ}=nx\), имеем
\[
\boxed{\Pi=nx-S}.
\]

Переплату также можно получить сложением всех начисленных процентов. Например,
для трёх периодов:
\[
\Pi
=
Sp
+
(Sr-x)p
+
\bigl((Sr-x)r-x\bigr)p.
\]

Оба способа дают одинаковый результат, однако формула
\[
\Pi=nx-S
\]
обычно значительно короче.

\section{Процент переплаты}

\term{Процент переплаты} показывает, на сколько процентов общая сумма выплат
больше первоначальной суммы кредита.

По определению
\[
\boxed{
\Pi_{\%}
=
\frac{\text{ОСВ}-S}{S}\cdot 100\%
}.
\]

Так как \(\text{ОСВ}=nx\), получаем
\[
\boxed{
\Pi_{\%}
=
\frac{nx-S}{S}\cdot 100\%
}.
\]

Для кредита на три периода:
\[
\Pi_{\%}
=
\left(
\frac{3x}{S}-1
\right)\cdot 100\%.
\]

Подставляя
\[
x=\frac{Sr^3}{r^2+r+1},
\]
получаем
\[
\Pi_{\%}
=
\left(
\frac{3Sr^3}{S(r^2+r+1)}-1
\right)\cdot 100\%.
\]

После сокращения \(S\):
\[
\boxed{
\Pi_{\%}
=
\left(
\frac{3r^3}{r^2+r+1}-1
\right)\cdot 100\%
}.
\]

\begin{checklist}
  \item Процент переплаты при фиксированных ставке и количестве периодов
  \important{не зависит от первоначальной суммы кредита}.

  \item Первоначальная сумма \(S\) сокращается в формуле процента переплаты.

  \item При увеличении процентной ставки переплата возрастает.

  \item При увеличении количества периодов общая сумма выплат, как правило,
  также возрастает.
\end{checklist}

\section{Каскад решения задачи}

Основной маршрут решения:
\[
\boxed{
\text{таблица}
\longrightarrow
x
\longrightarrow
\text{ОСВ}
\longrightarrow
\Pi
\longrightarrow
\Pi_{\%}
}
\]

\begin{checklist}
  \item Выпишите первоначальную сумму кредита \(S\).

  \item Выпишите количество периодов \(n\).

  \item Переведите процентную ставку:
  \[
  p=\frac{P}{100}.
  \]

  \item Введите коэффициент:
  \[
  r=1+p.
  \]

  \item Составьте таблицу изменения задолженности.

  \item Запишите остаток задолженности после последнего платежа.

  \item Приравняйте последний остаток к нулю.

  \item Выразите платёж \(x\).

  \item Найдите общую сумму выплат:
  \[
  \text{ОСВ}=nx.
  \]

  \item Найдите переплату:
  \[
  \Pi=nx-S.
  \]

  \item При необходимости найдите процент переплаты:
  \[
  \Pi_{\%}=\frac{nx-S}{S}\cdot100\%.
  \]

  \item Сопоставьте полученный результат с вопросом задачи.
\end{checklist}

\section{Типичные ошибки}

\begin{checklist}
  \item Использовать \(P\) вместо
  \(\displaystyle p=\frac{P}{100}\).

  \item Начислять проценты на первоначальную сумму \(S\) во все периоды,
  а не на текущий остаток задолженности.

  \item Вычитать платёж до начисления процентов, если условием установлен
  обратный порядок.

  \item Считать платежи разными, хотя схема является аннуитетной.

  \item Потерять одно из слагаемых при раскрытии вложенных скобок.

  \item Ошибиться в знаках:
  \[
  (Sr-x)r-x=Sr^2-xr-x.
  \]

  \item Записать после последнего платежа неравенство вместо условия
  \[
  D_n=0.
  \]

  \item Перепутать переплату
  \[
  \Pi=\text{ОСВ}-S
  \]
  с процентом переплаты
  \[
  \Pi_{\%}=\frac{\Pi}{S}\cdot100\%.
  \]

  \item Сократить выражения, не проверив, что сокращаемый множитель
  не равен нулю.

  \item Использовать формулу суммы геометрической прогрессии при \(r=1\)
  без отдельного рассмотрения случая \(p=0\).
\end{checklist}

\section{Минимальный пример}

Кредит равен
\[
S=331\,000\rub,
\]
ставка составляет \(10\%\) годовых, срок --- три года, платежи равные.

Переведём ставку:
\[
p=\frac{10}{100}=0{,}1,
\qquad
r=1+0{,}1=1{,}1=\frac{11}{10}.
\]

Платёж:
\begin{align*}
x
&=
\frac{Sr^3}{r^2+r+1}\\
&=
\frac{
331\,000\cdot\left(\frac{11}{10}\right)^3
}{
\left(\frac{11}{10}\right)^2+
\frac{11}{10}+1
}\\
&=
133\,100\rub.
\end{align*}

Общая сумма выплат:
\[
\text{ОСВ}=3x=399\,300\rub.
\]

Переплата:
\[
\Pi=399\,300-331\,000=68\,300\rub.
\]

Процент переплаты:
\[
\Pi_{\%}
=
\frac{68\,300}{331\,000}\cdot100\%
=
\frac{6830}{331}\%
\approx 20{,}63\%.
\]

% ============================================================
% Тренировочный блок
% ============================================================

\newpage

\section*{Тренировочный блок}

Во всех задачах предполагается, что проценты начисляются один раз в конце
каждого года перед внесением очередного равного платежа.

\subsection*{Средний уровень}

\begin{training}
  \item
  Кредит величиной \(S\) погашается тремя равными ежегодными платежами \(x\).
  Коэффициент увеличения задолженности равен \(r\).

  Запишите остатки задолженности \(D_1\), \(D_2\) и \(D_3\) после первого,
  второго и третьего платежей.

  \item
  Используя условие полного погашения кредита за три года, выведите формулу
  ежегодного платежа \(x\) через \(S\) и \(r\).

  \item
  Банк выдал кредит в размере \(331\,000\rub\) на три года под \(10\%\)
  годовых. Платежи в конце каждого года равны.

  Найдите:
  \begin{enumerate}[label=\alph*),leftmargin=2em]
    \item ежегодный платёж;
    \item общую сумму выплат;
    \item переплату;
    \item процент переплаты.
  \end{enumerate}
\end{training}

\subsection*{Уровень выше среднего}

\begin{training}[resume]
  \item
  Докажите, что после \(n\)-го равного платежа остаток задолженности равен
  \[
  D_n
  =
  Sr^n-x\left(r^{n-1}+r^{n-2}+\dots+r+1\right).
  \]
  Затем выразите \(x\).

  \item
  Первоначальная сумма кредита равна \(Y\), ставка в долях единицы равна
  \(t\), а
  \[
  k=1+t.
  \]
  Кредит погашается четырьмя равными платежами \(B\).

  Выведите формулу для \(B\) и разложите её знаменатель на множители.

  \item
  Кредит погашается двумя равными ежегодными платежами по
  \(121\,000\rub\). Годовая ставка составляет \(10\%\).

  Найдите:
  \begin{enumerate}[label=\alph*),leftmargin=2em]
    \item первоначальную сумму кредита;
    \item общую сумму выплат;
    \item переплату.
  \end{enumerate}

  \item
  Два заёмщика получили кредиты на одинаковое количество лет под одну и ту
  же процентную ставку, но первоначальные суммы кредитов различаются.

  Докажите, что процент переплаты у заёмщиков будет одинаковым.

  \item
  Кредит в размере \(331\,000\rub\) погашается тремя равными ежегодными
  платежами по \(133\,100\rub\).

  Найдите годовую процентную ставку, если она положительна.

  \item
  Кредит погашается за три года равными ежегодными платежами. Ставка
  составляет \(20\%\) годовых, а общая сумма выплат равна
  \(648\,000\rub\).

  Найдите:
  \begin{enumerate}[label=\alph*),leftmargin=2em]
    \item ежегодный платёж;
    \item первоначальную сумму кредита;
    \item переплату;
    \item процент переплаты.
  \end{enumerate}
\end{training}

\subsection*{Сложный уровень}

\begin{training}[resume]
  \item
  Для кредита на \(n\) периодов со ставкой \(p>0\) в долях единицы положите
  \(r=1+p\).

  Выведите формулу процента переплаты только через \(n\), \(p\) и \(r\).
  Докажите, что полученное выражение не зависит от первоначальной суммы
  кредита. Определите, к чему стремится процент переплаты при
  \(p\to0\).
\end{training}

% ============================================================
% Ответы
% ============================================================

\newpage

\section*{Ответы для самопроверки}

\begin{table}[ht]
  \centering
  \caption{Ответы к тренировочному блоку}
  \renewcommand{\arraystretch}{1.55}
  \small
  \begin{tabularx}{\linewidth}{
    @{}
    c
    >{\raggedright\arraybackslash}X
    @{}
  }
    \toprule
    \textbf{№} & \textbf{Ответ} \\
    \midrule

    1 &
    \[
    D_1=Sr-x,
    \]
    \[
    D_2=Sr^2-x(r+1),
    \]
    \[
    D_3=Sr^3-x(r^2+r+1).
    \]
    \\

    2 &
    \[
    \boxed{x=\frac{Sr^3}{r^2+r+1}}.
    \]
    \\

    3 &
    \[
    x=133\,100\rub,
    \]
    \[
    \text{ОСВ}=399\,300\rub,
    \qquad
    \Pi=68\,300\rub,
    \]
    \[
    \Pi_{\%}
    =
    \frac{6830}{331}\%
    \approx20{,}63\%.
    \]
    \\

    4 &
    \[
    D_n
    =
    Sr^n-x\sum_{j=0}^{n-1}r^j,
    \]
    \[
    \boxed{
    x=
    \frac{Sr^n}
    {r^{n-1}+r^{n-2}+\dots+r+1}
    }.
    \]
    При \(p>0\):
    \[
    \boxed{x=\frac{Spr^n}{r^n-1}}.
    \]
    \\

    5 &
    \[
    \boxed{
    B=\frac{Yk^4}{k^3+k^2+k+1}
    }.
    \]
    Поскольку
    \[
    k^3+k^2+k+1=(k+1)(k^2+1),
    \]
    то
    \[
    \boxed{
    B=\frac{Yk^4}{(k+1)(k^2+1)}
    }.
    \]
    \\

    6 &
    \[
    S=210\,000\rub,
    \qquad
    \text{ОСВ}=242\,000\rub,
    \qquad
    \Pi=32\,000\rub.
    \]
    \\

    7 &
    \[
    \Pi_{\%}
    =
    \left(
    \frac{nr^n}
    {r^{n-1}+r^{n-2}+\dots+r+1}
    -1
    \right)\cdot100\%.
    \]
    Первоначальная сумма кредита \(S\) в формуле отсутствует, поэтому
    процент переплаты одинаков.
    \\

    8 &
    \[
    \boxed{10\%}.
    \]
    \\

    9 &
    \[
    x=\frac{648\,000}{3}=216\,000\rub,
    \]
    \[
    S=455\,000\rub,
    \qquad
    \Pi=193\,000\rub,
    \]
    \[
    \Pi_{\%}
    =
    \frac{3860}{91}\%
    \approx42{,}42\%.
    \]
    \\

    10 &
    \[
    x=\frac{Spr^n}{r^n-1}.
    \]
    Поэтому
    \[
    \boxed{
    \Pi_{\%}
    =
    \left(
    \frac{npr^n}{r^n-1}-1
    \right)\cdot100\%
    }.
    \]
    Выражение не содержит \(S\), следовательно, процент переплаты не
    зависит от первоначальной суммы кредита. При \(p\to0\):
    \[
    x\to\frac{S}{n},
    \qquad
    \Pi_{\%}\to0\%.
    \]
    \\

    \bottomrule
  \end{tabularx}
\end{table}

\end{document}